\part{Understanding the system}

% A part divider is otherwise a page with four words on it. Giving each one a
% short orientation costs nothing and means no page in this book is bare.
\noindent\textbf{\large What this part covers}

\medskip
\noindent
Three chapters, and they are the ones to read slowly. The first explains what
the system does and the rules it enforces. The second shows the nine programs it
is built from and follows a single click all the way through them. The third
maps every directory, so that for the rest of the book you always know where a
file lives before you are shown what is in it.

\medskip
\noindent
Nothing after this point assumes you remember the details --- but it does assume
you have seen them once.

\chapter{What IncidentFlow actually is}
\label{ch:what}

\section{The problem}

When something breaks in production --- a payment service starts failing, a
database runs out of disk --- several things must happen at once. Someone has to
be told. Someone has to take charge. Everyone else needs to see what is being
done without interrupting the person doing it. And afterwards, someone has to be
able to reconstruct exactly what happened and when.

That last part is the hard one. Chat threads get edited. Memories disagree. The
engineer who fixed it at 3am writes the summary a week later.

\note{IncidentFlow's single organising idea is that \textbf{the timeline is
append-only}. Nothing written can be edited or deleted --- not by an
administrator, not by the person who wrote it, not by the API. If the record
says the incident was acknowledged at 02:14, that is what happened. Almost every
design decision in this book follows from that one commitment.}

\section{Who uses it}

Five roles, deliberately few. Each can do strictly more than the one above it.

\begin{center}
\begin{tabular}{@{}>{\raggedright\arraybackslash}p{4.2cm}>{\raggedright\arraybackslash}p{10.0cm}@{}}
\toprule
\textbf{Role} & \textbf{What it is for} \\
\midrule
\fpath{viewer} & Can read everything, change nothing. Support staff, managers, anyone who needs to know without being involved. \\
\fpath{reporter} & Can raise an incident, and comment. Cannot drive one. \\
\fpath{responder} & Works incidents: acknowledges, mitigates, resolves. \\
\fpath{incident_commander} & Owns an incident. Assigns people, changes severity, closes it out. \\
\fpath{administrator} & Manages the organisation: members, services, and the audit log. \\
\bottomrule
\end{tabular}
\end{center}

You will meet this table again in Chapter~\ref{ch:roles}, where the same list
appears as \emph{screenshots of one page seen by different people}. That is the
clearest way to understand a permission system: not as a table, but as buttons
that are missing.

\section{The shape of an incident}

Every incident has a \textbf{severity} --- how bad it is --- and a
\textbf{status} --- where it is in its life.

\subsection{Severity: how bad is it}

\begin{center}
\begin{tabular}{@{}llp{9.2cm}@{}}
\toprule
\textbf{Value} & \textbf{Name} & \textbf{Meaning} \\
\midrule
\fpath{sev1} & Critical & Everything is down, or money is being lost right now. Wakes people up. \\
\fpath{sev2} & High & A major feature is broken for many users. \\
\fpath{sev3} & Moderate & Degraded, or broken for a few. \\
\fpath{sev4} & Low & Cosmetic or minor; fix it in normal hours. \\
\bottomrule
\end{tabular}
\end{center}

\note{Severity is not decoration. \fpath{sev1} and \fpath{sev2} send email to
every responder in the organisation; \fpath{sev3} and \fpath{sev4} do not. One
enum value decides whether a hundred phones buzz at 3am --- which is why
changing severity is a permission a reporter does not have.}

\subsection{Status: where is it in its life}

\begin{center}
\begin{tabular}{@{}llp{9.0cm}@{}}
\toprule
\textbf{Value} & \textbf{Name} & \textbf{Meaning} \\
\midrule
\fpath{open} & Open & Reported. Nobody has picked it up yet. \\
\fpath{acknowledged} & Acknowledged & Someone is actively working on it. \\
\fpath{mitigated} & Mitigated & Users are no longer affected, but the underlying cause is still there. \\
\fpath{resolved} & Resolved & Actually fixed. \\
\fpath{closed} & Closed & Reviewed and finished. Terminal. \\
\bottomrule
\end{tabular}
\end{center}

\subsection{The transitions are a rule, not a suggestion}

You cannot move an incident from any status to any other. The permitted moves
are written down in exactly one place, and the API refuses everything else.

\where{api/app/Enums/IncidentStatus.php}
\begin{lstlisting}[language=PHP]
public function allowedTransitions(): array
{
    return match ($this) {
        // A trivial incident can go straight to resolved without ceremony.
        self::Open         => [self::Acknowledged, self::Resolved],
        self::Acknowledged => [self::Mitigated, self::Resolved],
        // Mitigation that does not hold sends the incident back to active work.
        self::Mitigated    => [self::Resolved, self::Acknowledged],
        // Resolved is reversible: an incident that recurs was never resolved.
        self::Resolved     => [self::Closed, self::Acknowledged],
        // Closed is terminal. Post-closure findings belong in the postmortem.
        self::Closed       => [],
    };
}
\end{lstlisting}

Read that carefully, because three of those five lines encode a real opinion
about how incidents work.

\begin{itemize}[leftmargin=1.4cm]
\item \textbf{Open can jump straight to Resolved.} Forcing every trivial
  incident through an acknowledgement step produces a system people work around.
\item \textbf{Mitigated can go \emph{back} to Acknowledged.} A workaround that
  stops holding is not progress, and the timeline should say so.
\item \textbf{Closed goes nowhere.} Once closed, the record is finished.
  Anything learned later belongs in the postmortem, which is a separate
  document --- not a quiet edit to history.
\end{itemize}

\warnbox{This rule lives in the \emph{enum}, not in the controller and not in
the React app. The frontend also hides buttons for moves that are not allowed
--- but that is a courtesy to the user, not a security control. Anyone can send
any HTTP request they like. If the rule lived only in the browser, it would not
be a rule at all. You will see this pattern throughout: \textbf{the UI prevents
mistakes, the API prevents abuse}, and only one of them is trusted.}

\section{Your first look}

Here is the sign-in screen --- the whole system's front door.

\shot{01-login-empty.png}{The sign-in screen at \fpath{http://localhost:8080}.
Chapter~\ref{ch:login} takes this page apart line by line.}

And here is what a signed-in incident commander sees: everything currently
happening.

\shot{05-incidents-list.png}{The incidents list, signed in as Daniel Okoye
(incident commander). Seventeen incidents, seeded so the system has a realistic
history to show.}

Do not worry about details yet --- Part~II gives each screen a full chapter. For
now notice only this: every coloured badge on that page is one of the enum
values from the tables above. The pictures and the code are the same thing.
